\Needspace{12\baselineskip}
\originalchapter{低速黏性流与薄层}
\originalsection{Stokes阻力和长程影响}
在 $Re\ll1$、惯性可忽略的尺度，$-\nabla p+\mu\nabla^2\mathbf u=0$，$\nabla\cdot\mathbf u=0$。流体以 $U$ 掠过半径 $a$ 的无滑移球，轴对称解用极角 $\theta$ 表示为
\begin{align*}
u_r&=U\cos\theta\left(1-\frac{3a}{2r}+\frac{a^3}{2r^3}\right),\\
u_\theta&=-U\sin\theta\left(1-\frac{3a}{4r}-\frac{a^3}{4r^3}\right),\\
p-p_\infty&=-\frac{3\mu Ua}{2r^2}\cos\theta.
\end{align*}
代入连续与Stokes方程可核验，$r=a$ 处两速度为零、远处为均匀流。将压力和黏性牵引投影到流向并对球面积分，各贡献 $2\pi\mu aU$ 和 $4\pi\mu aU$，总阻力 $6\pi\mu aU$。

二维无限圆柱的纯Stokes问题无法同时满足无滑移和无限远均匀流，这是Stokes悖论；远处即使局部速度扰动很小，累计惯性仍有作用。Oseen近似保留均匀背景对扰动的对流，修复远场匹配。这个例子说明小参数近似可能在空间远处不一致。本册保留这一物理边界，不给出未展开的完整圆柱匹配系数。

\originalsection{润滑近似}
薄隙高度 $h$ 远小于水平尺度 $L$，竖向速度小于水平速度，压力跨隙近似均匀。还需检查惯性相对于跨隙黏性足够小。在水平薄隙中 $\mu u_{zz}=p_x$，两壁静止给 $u=p_x z(z-h)/(2\mu)$，单位宽通量 $q=-h^3p_x/(12\mu)$。若上壁无切向应力、下壁无滑移，通量改为 $-h^3p_x/(3\mu)$，因子四来自边界条件改变。

重力作用下倾斜平面自由层，设 $z$ 垂直底面，坡角 $\alpha$，自由表面压力为常量，局部驱动力 $G=\rho g(\sin\alpha-\cos\alpha\,h_x)$。速度 $u=Gz(2h-z)/(2\mu)$，通量 $q=Gh^3/(3\mu)$，质量守恒给 $h_t+q_x=0$。局部层厚为恒定时恢复抛物线剖面；空间变化时曲率或长坡度产生展平作用。

\originalsection{重力流的相似尺度}
水平面的黏性重力流满足 $h_t=(g/3\nu)(h^3h_x)_x$。一维固定体积每单位宽度 $V\sim hR$，由时间尺度 $t\sim\nu R^2/(g h^3)$ 得 $R\sim t^{1/5}$、$h\sim t^{-1/5}$。轴对称固定体积 $V\sim hR^2$ 则给 $R\sim t^{1/8}$、$h\sim t^{-1/4}$。几何改变质量约束，因而改变相似指数；不能在轴对称熔岩题中照搬一维指数。

Bingham流体还需超过屈服应力才能剪切。若底部应力低于屈服应力，薄层可以完全静止；超过后底部形成剪切层，上部可作近似塞流。Newton黏性公式并不能仅把一个“有效黏度”随意调大就还原屈服阈值。气溶胶凝并、选择性抽取等源课主题有其他物理机制，未在本版建立其完整专门模型。

\originalsection{平板与圆管：黏性和压力驱动的基准}
两平板间高度 $h$，稳态沿 $x$、$u=u(y)$、$p_x$ 常数，NS化为 $\mu u''=p_x$。两壁静止得 $u=-p_x y(h-y)/(2\mu)$，单位宽流量 $q=-h^3p_x/(12\mu)$；一壁以 $U$ 移动则加 $Uy/h$。圆管半径 $a$ 对应 $u=-p_x(a^2-r^2)/(4\mu)$、$Q=-\pi a^4p_x/(8\mu)$。厚度三次/半径四次依赖解释为什么薄间隙误差显著放大流量误差。

入口未发展、惯性、非圆截面或非Newton性质会改变这些式子。壁滑移会改变边界而非仅改一个黏度。以它们校准黏度时，先核查几何和压力落差是否主要发生在测试段。

\originalsection{润滑方程的完整推导}
薄层 $h/L=\varepsilon\ll1$，纵向速度 $U$、垂向速度由连续式为 $O(\varepsilon U)$。若 $Re\,\varepsilon^2\ll1$，横向黏性 $\mu u_{yy}$ 主导纵向黏性与惯性，垂向动量给 $p_y\simeq-\rho g$。刚性间隙的局部流量 $q=-h^3p_x/(12\mu)+(U_1+U_2)h/2$。质量守恒 $h_t+q_x=0$ 得Reynolds润滑方程。

自由表面在上端切向应力零，底部无滑移，则压力驱动系数改为 $1/(3\mu)$。平地重力流 $p_x=\rho gh_x$，因此 $h_t=(g/(3\nu))\partial_x(h^3h_x)$。若加表面张力，近小坡度表压为 $-\gamma_sh_{xx}$，通量含 $+\gamma_sh^3h_{xxx}/(3\mu)$。相同的连续式用不同应力条件产生不同系数。

黏附两板的外缘压设大气，打开间隙需要负表压，源题可由润滑压力积分求开口速度；若压力低于气穴阈值，模型失效。数学上趋无限大的吸力并非真实流体可以无限承受张力。

\originalsection{Stokes球流与阻力}
小Re球外稳态满足 $-\nabla p+\mu\nabla^2\mathbf u=0,\nabla\cdot\mathbf u=0$，球面无滑移、远场均速 $U\hat z$。轴对称解（球半径 $a$）为
\begin{align*}
u_r&=U\cos\theta\left(1-\frac{3a}{2r}+\frac{a^3}{2r^3}\right),\\
u_\theta&=-U\sin\theta\left(1-\frac{3a}{4r}-\frac{a^3}{4r^3}\right),\\
p-p_\infty&=-\frac{3\mu Ua}{2r^2}\cos\theta.
\end{align*}
直接检查 $r=a$ 两速度零、$r\to\infty$ 恢复均匀流。把球面压力和剪切牵引沿流向积分，压力贡献 $2\pi\mu aU$、剪切贡献 $4\pi\mu aU$，总阻力 $6\pi\mu aU$。只积分压力会漏三分之二。

粒子终端沉降满足 $6\pi\mu aU_t=4\pi a^3(\rho_p-\rho_f)g/3$，故 $U_t=2a^2(\rho_p-\rho_f)g/(9\mu)$。再以该速度算Re，检查最初小Re假设；粒子离壁或邻粒太近也改变阻力。

\originalsection{Oseen远场和Stokes悖论}
即使球附近Re小，远场扰动导数越来越小，背景速度的对流 $U\partial_z\mathbf u'$ 与黏性 $\nu\nabla^2\mathbf u'$ 的比约 $Ur/\nu$，在 $r\sim a/Re_a$ 变成一阶。故小Re展开并非全空间一致，不能在无穷远都删除惯性。Oseen线性化保留背景对流，得到 $\rho U\partial_z\mathbf u'=-\nabla p'+\mu\nabla^2\mathbf u'$。

球的内Stokes解与外Oseen解在 $a\ll r\ll a/Re_a$ 重叠匹配，第一阶阻力为 $6\pi\mu aU[1+3Re_a/8+\cdots]$，其中 $Re_a=Ua/\nu$；若以直径定义Re，系数为 $3Re_d/16$。本节给匹配的尺度、方程和已知首阶结果，未把未展开的完整外场系数推导说成已证明。

二维无限圆柱的Stokes解无法同时满足静止无滑移与无穷远均速（Stokes悖论），说明远场惯性即使小也不能完全略去。这不是NS没有物理解，而是近似丢了决定远场的项。源课Oseen主题的核心就在这个奇异极限，而非只背一个Re数。

\originalsection{气溶胶的惯性、扩散与连续介质限制}
小球的速度弛豫时间 $\tau_p=m_p/(6\pi\mu a)=2\rho_pa^2/(9\mu)$。流场变化时间 $T_f$ 下 $St=\tau_p/T_f\ll1$ 时粒子近随流，$St$大则偏离流线产生惯性撞击。Brown扩散由 $D_p=\kb T/(6\pi\mu a)$ 给出，小粒子扩散更快而沉降更慢，形成气溶胶输运机制的竞争。

气体平均自由程 $\lambda_g$ 与半径接近时，Knudsen数 $Kn=\lambda_g/a$ 不小，连续无滑移Stokes阻力需要滑移修正或分子模型。修正参数应来自相应资料/标定；本文不捏造某种粒子的实测经验系数。过滤/采样装置还需将扩散、惯性、重力与几何接受联合考虑。
